% GENERATED FILE. Edit canonical lessons, then run npm run generate:books.
% canonical-source-hash: fnv1a64:0821c151dbe08b1d
% canonical-lessons: GU-C201-sudharvu, GU-C201-moci, GU-C201-dhobi, GU-C201-hajam, GU-C201-suthar

\chapter{Getting things mended}
\label{ch:gu-getting-things-mended}
\hlchaptermodality{201}

\begin{chapteropening}
\textbf{By the end of this chapter:} \emph{I can name the people who mend, wash and cut things for me.}

Complete the last lesson of chapter 201: I can name the people who mend, wash and cut things for me.
\end{chapteropening}

\section[\texorpdfstring{\gu{સુધારવું} (sudhārvũ)}{સુધારવું (sudhārvũ)}]{\gu{સુધારવું} (sudhārvũ) — to mend, to repair, to correct}

\label{lesson:GU-C201-sudharvu}



\begin{quote}
Before the new one: say the Gujarati for to put in, then the Gujarati for to register.
\end{quote}

\subsection*{You'll want to know: \gu{સુધારવું}}
\textbf{\gu{સુધારવું}} — \emph{sudhārvũ} — \textquotedblleft{}to mend, to repair, to correct\textquotedblright{}.

\begin{quote}
\textbf{\gu{તે} \gu{જૂની} \gu{સાઈકલ} \gu{સુધારે} \gu{છે}.} — \emph{te jūnī saikal sudhāre che} — \textquotedblleft{}He repairs the old bicycle.\textquotedblright{}
\end{quote}

\begin{cousinweb}
\textbf{\gu{સુધારવું}} (\emph{sudhārvũ}) is a tadbhava. CDIAL traces it to \emph{*śuddhakāra}, \textquotedblleft{}making clean\textquotedblright{}, built from Sanskrit \emph{śuddhá}, \textquotedblleft{}clean\textquotedblright{}, and \emph{kāra}, \textquotedblleft{}making\textquotedblright{} (CDIAL 12521).
\end{cousinweb}

\begin{grammarlens}[title={What you've built}]
One more word for everyday life at A2.
\end{grammarlens}

\subsection*{Your turn}
\begin{itemize}
\raggedright
  \item \emph{Say it:} \emph{sudhārvũ}
  \item \emph{Say it:} \emph{sudhārvũ}, once more
  \item \emph{Say it:} \emph{te jūnī saikal sudhāre che} — He repairs the old bicycle.
  \item \emph{Recall:} say the Gujarati for to put in, then the Gujarati for to register, then say \emph{sudhārvũ} again
\end{itemize}

\begin{tcolorbox}[breakable,colback=teal!4,colframe=teal!35!black,title={Before you move on}]
What does \gu{સુધારવું} mean? (\textquotedblleft{}To mend, to repair, to correct\textquotedblright{}.) Say it once more.
\end{tcolorbox}
\section[\texorpdfstring{\gu{મોચી} (mocī)}{મોચી (mocī)}]{\gu{મોચી} (mocī) — cobbler}

\label{lesson:GU-C201-moci}



\begin{quote}
Before the new one: say the Gujarati for to register, then the Gujarati for to mend.
\end{quote}

\subsection*{You'll want to know: \gu{મોચી}}
\textbf{\gu{મોચી}} — \emph{mocī} — \textquotedblleft{}cobbler\textquotedblright{}.

\begin{quote}
\textbf{\gu{મોચીએ} \gu{મારા} \gu{બૂટ} \gu{સુધાર્યા}.} — \emph{mocīe mārā būṭ sudhāryā} — \textquotedblleft{}The cobbler mended my shoes.\textquotedblright{}
\end{quote}

\begin{cousinweb}
\textbf{\gu{મોચી}} (\emph{mocī}) comes from Sanskrit \emph{mōcika}, \textquotedblleft{}tanner and cobbler\textquotedblright{}, built on \emph{mōca}, \textquotedblleft{}shoe\textquotedblright{}, which CDIAL takes from Iranian, as in Pahlavi \emph{mōč}, \textquotedblleft{}shoe\textquotedblright{}; CDIAL glosses the Gujarati word \textquotedblleft{}cobbler\textquotedblright{} (CDIAL 10349).
\end{cousinweb}

\begin{grammarlens}[title={What you've built}]
One more word for everyday life at A2.
\end{grammarlens}

\subsection*{Your turn}
\begin{itemize}
\raggedright
  \item \emph{Say it:} \emph{mocī}
  \item \emph{Say it:} \emph{mocī}, once more
  \item \emph{Say it:} \emph{mocīe mārā būṭ sudhāryā} — The cobbler mended my shoes.
  \item \emph{Recall:} say the Gujarati for to register, then the Gujarati for to mend, then say \emph{mocī} again
\end{itemize}

\begin{tcolorbox}[breakable,colback=teal!4,colframe=teal!35!black,title={Before you move on}]
What does \gu{મોચી} mean? (\textquotedblleft{}Cobbler\textquotedblright{}.) Say it once more.
\end{tcolorbox}
\section[\texorpdfstring{\gu{ધોબી} (dhobī)}{ધોબી (dhobī)}]{\gu{ધોબી} (dhobī) — washerman}

\label{lesson:GU-C201-dhobi}



\begin{quote}
Before the new one: say the Gujarati for to mend, then the Gujarati for cobbler.
\end{quote}

\subsection*{You'll want to know: \gu{ધોબી}}
\textbf{\gu{ધોબી}} — \emph{dhobī} — \textquotedblleft{}washerman\textquotedblright{}.

\begin{quote}
\textbf{\gu{ધોબી} \gu{ગલીના} \gu{નાકે} \gu{છે}.} — \emph{dhobī galīnā nāke che} — \textquotedblleft{}The washerman is at the end of the lane.\textquotedblright{}
\end{quote}

\begin{cousinweb}
\textbf{\gu{ધોબી}} (\emph{dhobī}) comes from \emph{*dhauvati}, \textquotedblleft{}washes\textquotedblright{}, a reconstructed form beside Sanskrit \emph{dhautá}, \textquotedblleft{}washed\textquotedblright{}. CDIAL says the forms with \emph{-b-} began in the east, and that Hindi \emph{dhobī} spread from there into the other languages, Gujarati among them, so the Gujarati word is a loan from Hindi (CDIAL 6886).
\end{cousinweb}

\begin{grammarlens}[title={What you've built}]
One more word for everyday life at A2.
\end{grammarlens}

\subsection*{Your turn}
\begin{itemize}
\raggedright
  \item \emph{Say it:} \emph{dhobī}
  \item \emph{Say it:} \emph{dhobī}, once more
  \item \emph{Say it:} \emph{dhobī galīnā nāke che} — The washerman is at the end of the lane.
  \item \emph{Recall:} say the Gujarati for to mend, then the Gujarati for cobbler, then say \emph{dhobī} again
\end{itemize}

\begin{tcolorbox}[breakable,colback=teal!4,colframe=teal!35!black,title={Before you move on}]
What does \gu{ધોબી} mean? (\textquotedblleft{}Washerman\textquotedblright{}.) Say it once more.
\end{tcolorbox}
\section[\texorpdfstring{\gu{હજામ} (hajām)}{હજામ (hajām)}]{\gu{હજામ} (hajām) — barber}

\label{lesson:GU-C201-hajam}



\begin{quote}
Before the new one: say the Gujarati for cobbler, then the Gujarati for washerman.
\end{quote}

\subsection*{You'll want to know: \gu{હજામ}}
\textbf{\gu{હજામ}} — \emph{hajām} — \textquotedblleft{}barber\textquotedblright{}.

\begin{quote}
\textbf{\gu{હજામની} \gu{દુકાન} \gu{ચોકમાં} \gu{છે}.} — \emph{hajāmnī dukān cokmā̃ che} — \textquotedblleft{}The barber's shop is in the square.\textquotedblright{}
\end{quote}

\begin{cousinweb}
\textbf{\gu{હજામ}} (\emph{hajām}) is a loan from Arabic \emph{ḥajjām}, \textquotedblleft{}cupper\textquotedblright{}, one who let blood as a cure, which reached Gujarati through Persian. Its Arabic root is \emph{ḥ-j-m} (Lane).
\end{cousinweb}

\begin{grammarlens}[title={What you've built}]
One more word for everyday life at A2.
\end{grammarlens}

\subsection*{Your turn}
\begin{itemize}
\raggedright
  \item \emph{Say it:} \emph{hajām}
  \item \emph{Say it:} \emph{hajām}, once more
  \item \emph{Say it:} \emph{hajāmnī dukān cokmā̃ che} — The barber's shop is in the square.
  \item \emph{Recall:} say the Gujarati for cobbler, then the Gujarati for washerman, then say \emph{hajām} again
\end{itemize}

\begin{tcolorbox}[breakable,colback=teal!4,colframe=teal!35!black,title={Before you move on}]
What does \gu{હજામ} mean? (\textquotedblleft{}Barber\textquotedblright{}.) Say it once more.
\end{tcolorbox}
\section[\texorpdfstring{\gu{સુથાર} (suthār)}{સુથાર (suthār)}]{\gu{સુથાર} (suthār) — carpenter}

\label{lesson:GU-C201-suthar}



\begin{quote}
Before the new one: say the Gujarati for washerman, then the Gujarati for barber.
\end{quote}

\subsection*{You'll want to know: \gu{સુથાર}}
\textbf{\gu{સુથાર}} — \emph{suthār} — \textquotedblleft{}carpenter\textquotedblright{}.

\begin{quote}
\textbf{\gu{સુથારે} \gu{નવું} \gu{બારણું} \gu{બનાવ્યું}.} — \emph{suthāre navũ bārṇũ banāvyũ} — \textquotedblleft{}The carpenter made a new door.\textquotedblright{}
\end{quote}

\begin{cousinweb}
\textbf{\gu{સુથાર}} (\emph{suthār}) is a tadbhava. It comes from Sanskrit \emph{sūtradhāra}, \textquotedblleft{}carpenter\textquotedblright{}, built from \emph{sū́tra}, \textquotedblleft{}thread, line\textquotedblright{}, and \emph{dhāra}, \textquotedblleft{}holding\textquotedblright{}, and CDIAL glosses the Gujarati word \textquotedblleft{}carpenter\textquotedblright{} (CDIAL 13563).
\end{cousinweb}

\begin{grammarlens}[title={What you've built}]
One more word for everyday life at A2.
\end{grammarlens}

\subsection*{Your turn}
\begin{itemize}
\raggedright
  \item \emph{Say it:} \emph{suthār}
  \item \emph{Say it:} \emph{suthār}, once more
  \item \emph{Say it:} \emph{suthāre navũ bārṇũ banāvyũ} — The carpenter made a new door.
  \item \emph{Recall:} say the Gujarati for washerman, then the Gujarati for barber, then say \emph{suthār} again
\end{itemize}

\begin{tcolorbox}[breakable,colback=teal!4,colframe=teal!35!black,title={Before you move on}]
What does \gu{સુથાર} mean? (\textquotedblleft{}Carpenter\textquotedblright{}.) Say it once more.
\end{tcolorbox}
